\section{What the models rely on}
\label{explain}

\subsection{TF-IDF logistic regression: global coefficients}

\textbf{Setup.} TF-IDF LR refitted on the seed-0 training rows with its tuned settings ($C = 10$,
30,000 terms, 1- and 2-grams). Table~\ref{tab:terms} lists the ten largest coefficients on each side.

\begin{table}[tb]
\centering
\small
\begin{tabular}{lrlr}
\toprule
\multicolumn{2}{c}{\textbf{Towards fraud}} & \multicolumn{2}{c}{\textbf{Towards real}} \\
\textbf{Feature} & \textbf{Coef.} & \textbf{Feature} & \textbf{Coef.} \\
\midrule
\tabinput{includes/tab_terms}
\bottomrule
\end{tabular}
\caption{Largest coefficients of TF-IDF LR (seed 0). Plain entries are text terms; ``industry'' and
``country'' entries are one-hot metadata columns.}
\label{tab:terms}
\end{table}

\textbf{Money words.} \emph{money}, \emph{income} and \emph{000} (from amounts such as 50,000) push
towards fraud, in line with the money-word flag of Section~\ref{data}.

\textbf{Poster templates.} The phrase \emph{apply using} (the stop word in ``apply using the link''
is dropped) appears in 41 ads, 35 of them fraudulent; \emph{accion} is an employer name found in 26
ads, 16 of them fraudulent. These are the wording of a few repeated posters, not general scam
language. The grouped split keeps ads with the same description together, but a poster's
different descriptions (41 groups for \emph{apply using}) still reach both sides of the split.

\textbf{Signals of a real employer.} Established-company words (\emph{growing}, \emph{companies}),
specific industries (insurance, internet, restaurants) and the countries Greece and the Philippines
push towards real. The country entries are the collection artefact flagged in Section~\ref{data}.

\subsection{DistilBERT: Integrated Gradients on single ads}

\textbf{Setup.} Ten seed-0 test ads, chosen by score as in Section~\ref{method}. Table~\ref{tab:ig}
lists, per ad, the words with the largest push towards fraud, after merging word pieces and dropping
stop words and punctuation (\$ and numbers are kept).

\textbf{Check.} Each ad's attributions add up to $+0.50$ to $+0.53$ for the flagged ads and
$-0.48$ for the missed ones, which is the change in fraud probability from the all-\texttt{[PAD]}
baseline (about 0.5) to the ad itself. Integrated Gradients must satisfy this completeness property,
so the attributions are numerically sound.

\begin{table}[tb]
\centering
\footnotesize
\setlength{\tabcolsep}{4pt}
\begin{tabular}{llrr>{\raggedright\arraybackslash}p{0.36\linewidth}}
\toprule
\textbf{Kind} & \textbf{Title} & \textbf{Score} & \textbf{Sum} & \textbf{Top words towards fraud} \\
\midrule
\tabinput{includes/tab_ig}
\bottomrule
\end{tabular}
\caption{Integrated Gradients for DistilBERT (seed 0). Score is the fraud probability, Sum the sum
of all token attributions. A word can be listed twice when it occurs twice in the ad.}
\label{tab:ig}
\end{table}

\textbf{Caught scams.} The two clearest scams (part-time and weekend cash jobs) are pushed by pay
amounts and money words: \$, 500, 400, 450, \emph{earn}. These are the same cues as the money-word
flag and the linear model's top terms.

\textbf{False alarms.} The three real ads DistilBERT is most sure about are part-time, data-entry
and payroll-clerk ads that quote a pay rate (\$20/hr). They share the surface features of the
scams, so a model that reads only the text cannot separate them.

\textbf{Missed frauds.} The three missed frauds (\emph{DATA ENTRY}, \emph{Front-End Developer},
\emph{Hotel Workers Needed}) read like ordinary job descriptions; their attributions sum to $-0.48$,
so the text as a whole pushes towards real and no single word stands out.

\textbf{Artefact.} One caught fraud is pushed by a token like \texttt{ddb08035...}: the hash that
EMSCAD inserted in place of a real URL. A hash identifies one specific link, so the model has
partly memorised a poster rather than a pattern.

\subsection{Error analysis}

\textbf{DistilBERT versus TF-IDF LR.} At their own thresholds on the 171 seed-0 test frauds, both
catch 109, only DistilBERT catches 12, only TF-IDF LR catches 15, and 35 get past both. The
Spearman correlation between their test scores is only 0.39. Two models with similar PR-AUC
that rank ads this differently make different mistakes, which is the condition for stacking to help
(Section~\ref{stacking}). DistilBERT raises more false alarms at its threshold (28 against 19).

\textbf{Frauds no model catches.} 18 seed-0 test frauds are missed by every supervised model at its
own threshold:
\begin{itemize}[nosep]
  \item 11 IT and engineering roles written like ordinary postings (business analyst, SAP
        engineer, integration engineer, process engineer), four of them in one shared
        ``Hiring for ... Position: ... Job Type: Full Time'' layout;
  \item one call-centre ad present three times, all in the test split because the split is grouped;
  \item two hotel jobs in London and one vacancy that names a real oil and gas company;
  \item one customer-assistant ad.
\end{itemize}
Only 3 of the 18 contain money words, against 43\% of all seed-0 test frauds, and several borrow
the name of a real employer. Their text barely differs from real ads; catching them would need information outside the ad, such
as the poster's account history.
